\section{Related Work}\label{sec:related-work}

The field of ACR has seen considerable research since the seminal work of~cite{FujishimaACR} in 1999. Below, we provide a brief overview of the field including the models that are commonly used and the methods used to improve performance. We then highlight the use of generative models and the use of frames and beats in ACR and other music information retrieval tasks in order to motivate the methods we apply in this work.

\textbf{Models:} Increasingly complex models have been applied, following the trends in other fields. This started with CNNs~\cite{RethinkingChordRecognition}, then these were combined with recurrent neural networks (RNNs)~\cite{firstrnn, ACRCNNRNN1,ACRLargeVocab1,StructuredTraining} with a CNN performing feature extraction from a spectrogram and an RNN sharing information across frames. We refer to these models as CRNNs. More recently, transformers have been applied in place of the RNN~\cite{MelodyTranscriptionViaGenerativePreTraining, HarmonyTransformer, AttendToChords,CurriculumLearning,BTC,discordharmony,chordformer}. \cite{FourTimelyInsights} discuss the `glass ceiling' in ACR performance that is still largely present today. The `BTC' model introduced by~\cite{BTC} has become a common baseline despite the lack of performance improvement in the original publication. The more recent `Chordformer'~\cite{chordformer} shows modest improvement of up to 2\% improvement in accuracy. Instead, for reasons of computational constraint and the wish to run many experiments, we choose to use the CRNN from~\cite{StructuredTraining} as our baseline model. This remains competitive with state-of-the-art and we hope that performance improvements on this model are representative of the possible improvement with more complex models.

\textbf{Improvement Methods:} \emph{Ouput smoothing} — Frame-by-frame chord prediction can create unstable outputs with many short chord transitions. To this end, the use of hidden mark models (HMMs)~\cite{BalanceRandomForestACR} and conditional random fields (CRFs)~\cite{ACRLargeVocab1} have been used as a post-processing step for output smoothing. \emph{Pitch Augmentation} — First introduced by~\cite{StructuredTraining}, models are trained with pitch alterations as a form of data augmentation to encourage root-invariance in the model and to increase the model's ability to generalise. \emph{Structured Decomposition} — It is now common to decompose the chord vocabulary into root, quality and extensions and to have the model predict these separately~\cite{StructuredTraining, ACRLargeVocab1}. Such decomposition is thought to increase performance by around 1\%, with greater gains at recognising the root, third and seventh. \emph{Weighting the loss} — To address the long tail of the chord label distribution, it is common to increase the weight of the loss on rare chord classes~\cite{ACRLargeVocab1}. We perform extensive ablations on four of these improvement methods: output smoothing, weighting the loss function, structured decomposition, pitch augmentation. 

Other methods have been attempted but are not explored in this work as they are less common or computationally too expensive for us. For example, some have suggested embracing the subjective nature of chord annotation~\cite{FourTimelyInsights,20YearsofACR}. This has been attempted by~\cite{discordharmony} who use label-smoothing informed by a consonance-based distance metric between chord labels.

\textbf{Generative Models:} Generative models have been transformative in other domains such as natural language processing and computer vision. In such domains, large generative models have become the dominant method for performing all classification and generation tasks. It is common for classification to be performed on embeddings produced by large, pre-trained, generative models, and to use generative models to create large synthetic datasets for further training. In music, generative models such as Jukebox~\cite{Jukebox} and MusicLM~\cite{MusicLM} have transformed the ability to quickly create some styles of music. They have been employed for music transcription of melody and chords~\cite{MelodyTranscriptionViaGenerativePreTraining} but extensive analysis for ACR using the dominant dataset is missing. We propose to follow the example set byu other domains and use generative models for ACR in two ways: (1) using features extracted from a generative model as input to a CRNN and (2) using a chord-conditioned generative model to create synthetic data to train on. We use generative features extracted from MusicGen~\cite{MusicGen} and use a MusiConGen~\cite{MusiConGen}, a chord-conditioned generative model, to generate synthetic training data for ACR.

\textbf{Frames and Beats:} Chords exist in time. ACR is most commonly performed on a constant-Q transform (CQT) of the audio signal, with the frame length determining the temporal resolution. A variety of frame lengths have been used to perform ACR, ranging from 512~\cite{ACRLargeVocab1} to 4096~\cite{StructuredTraining}. By contrast,~\cite{MelodyTranscriptionViaGenerativePreTraining} perform melody and chord transcription on a Mel-spectrogram~\cite{MelScale} by resampling generative features onto a frame length that is song-dependent and beat-aware. They use \texttt{madmom}~\cite{madmom} for beat estimation and resample their features to a grid of 1/16th notes. We propose to use beat detection models from \cite{madmom} to first identify beats in the song, and to investigate the relationship between varying the granularity of the beat grid and the performance of the ACR model. We also perform ablations on the CQT frame length and the type of spectrogram used as input to ensure that current best practices are reasonable.